\section{The exact iPRA backend}

This section fixes the object theory used throughout.  The details matter,
because several tempting shortcuts are invalid in so weak and syntactic a
setting.

\subsection{Primitive-recursive terms}

For each arity \(k\), the repository datatype \(\mathrm{PR}(k)\) is generated
by zero, successor, projections, composition, and primitive recursion.  A term
\(\mathrm{PTm}(n)\) over \(n\) free variables is either a variable or an
application of a fixed PR symbol.  Simultaneous substitution is defined by
structural recursion.

The theory \(\PRAAx\) contains the defining equations for every PR constructor,
successor separation and injectivity, and induction for quantifier-free
formulas.  Derivability
\[
   \Gamma\objectvdash\varphi
\]
is the inductive natural-deduction relation \leanname{Prf PRAAx Gamma phi}.
The logical rules are intuitionistic.  Equality is reflexive and admits
simultaneous Leibniz substitution.

\subsection{Structural facts used by the framework}

Four checked rules form the logical spine of \(\FQNF\).

\begin{enumerate}
\item \textbf{Substitution.}
A derivation may be specialized by arbitrary object terms.

\item \textbf{Weakening.}
A derivation remains valid under additional assumptions.

\item \textbf{Simultaneous cut.}
If \(\Delta\objectvdash\varphi\) and every member of \(\Delta\) is derivable
from \(\Gamma\), then \(\Gamma\objectvdash\varphi\).  This is
\leanname{Prf.trans}.

\item \textbf{Quantifier-free induction.}
For a QF formula \(I(z,\vec x)\), object proofs of the base and successor step
yield \(\forall z\,I(z,\vec x)\).  Parameters \(\vec x\) are allowed and may
include quantified hypotheses in the surrounding context; only the induction
matrix must be QF.
\end{enumerate}

The last point is often misunderstood.  The final theorem may be quantified or
contain implications and existentials.  What must remain quantifier-free is the
formula propagated by induction.  A terminal implication can extract an
arbitrary conclusion after the finite loop.

\subsection{Arithmetic components already available}

The checked derivation library supplies almost all low-level ingredients needed
for the normal form.

\begin{itemize}
\item natural addition, multiplication, predecessor, truncated subtraction,
and order;
\item finite sums, products, powers, and factorials;
\item a Boolean equality test \(\mathsf{eqb}\);
\item pairing with proved projections and simultaneous recursion;
\item quotient and remainder with Euclidean existence and uniqueness;
\item bounded least search with correctness and minimality;
\item bounded-sum reflection for finite universal conditions;
\item generic recursion-invariant patterns and matrix/ring tools.
\end{itemize}

The framework therefore does not propose a second arithmetic foundation.  It
organizes these ingredients into a reusable certificate compiler.

\subsection{Three forbidden inferences}

\begin{warning}[Host truth is not object derivability]
A Lean proof that a natural-number function has a property does not create a
\(\Prf\,\PRAAx\) derivation of the encoded formula.  Every semantic lemma used
by a certificate checker must have a reflected object proof.
\end{warning}

\begin{warning}[A Lean datatype is not a universal PR evaluator]
Although \(\mathrm{PR}\) is an inductive datatype in Lean, there is no total
primitive-recursive evaluator uniformly interpreting arbitrary PR syntax on
all inputs.  \(\FQNF\) specializes every program before internalization.
\end{warning}

\begin{warning}[Finiteness is not a cutoff]
A proposition that a set is finite does not supply a numerical bound below
which a primitive-recursive search can stop.  The bound is a separate producer
field and may contain the deepest mathematics in the application.
\end{warning}

These warnings are not bureaucratic.  They exactly locate the failures in
several ambitious proof packets.
